\chapter{Algebra}
\label{sec:algebra}

The algebra module is organized around two basic classes of structures:
groups and rings. Groups provide a single operation with inverses, while rings
combine an additive group structure with a compatible multiplication.

\section{Groups}

A group is a set $G$ equipped with a binary operation
\[
  \ast : G \times G \to G
\]
such that the operation is associative, has an identity element $e$, and every
element $a \in G$ has an inverse $a^{-1}$. In PrimePie this abstract structure
is represented by the common operations
\[
  \texttt{identity}, \qquad
  \texttt{inverse}, \qquad
  \texttt{operate}.
\]

\subsection*{Powers in a Group}

For a group element $a$ and an integer $n$, the power $a^n$ means repeated use
of the group operation:
\[
a^n =
\begin{cases}
e, & n = 0,\\
\underbrace{a \ast a \ast \cdots \ast a}_{n \text{ times}}, & n > 0,\\
(a^{-1})^{-n}, & n < 0.
\end{cases}
\]
This definition applies to any group, not only to multiplicative groups. For an
additive group, $a^n$ means repeated addition of $a$.

The algorithm used for group powers is binary exponentiation. For positive
$n$, it computes powers in $\mathcal{O}(\log n)$ group operations by repeatedly
squaring:
\[
a^n =
\begin{cases}
e, & n = 0,\\
(a^{n/2})^2, & n \text{ even},\\
a \ast (a^{(n-1)/2})^2, & n \text{ odd}.
\end{cases}
\]

\subsection*{Concrete Group Families}

The current algebra layer includes modular additive groups, modular
multiplicative groups of units, and direct sums of groups.

\begin{itemize}
  \item The additive group $\mathbb{Z}/n\mathbb{Z}$ uses addition modulo $n$.
  \item The multiplicative group $(\mathbb{Z}/n\mathbb{Z})^\times$ consists of
        residue classes coprime to $n$, with multiplication modulo $n$.
  \item Direct sums combine several groups componentwise:
        \[
          (a_1,\ldots,a_k) \ast (b_1,\ldots,b_k)
          =
          (a_1 \ast_1 b_1,\ldots,a_k \ast_k b_k).
        \]
\end{itemize}

\section{Rings}

A ring is a set $R$ equipped with addition and multiplication. The additive
structure is an abelian group:
\[
  (R,+,0),
\]
and multiplication is associative and distributes over addition:
\[
a(b+c)=ab+ac,
\qquad
(a+b)c=ac+bc.
\]
The ring structures in PrimePie expose the canonical operations
\[
\texttt{zero}, \qquad
\texttt{one}, \qquad
\texttt{add}, \qquad
\texttt{neg}, \qquad
\texttt{mul}.
\]

\subsection*{Additive and Multiplicative Powers}

Since a ring is an additive group, it has an additive power operation inherited
from the group structure. Thus, in a ring,
\[
  \texttt{power}(a,n)
\]
means repeated addition of $a$.

Multiplicative exponentiation is a separate operation:
\[
  \texttt{mpower}(a,n),
\]
which computes
\[
a^n =
\underbrace{a \cdot a \cdots a}_{n \text{ times}}
\]
using binary exponentiation. In a general ring this operation is defined for
$n \ge 0$.

\subsection*{The Ring of Integers \texorpdfstring{$\mathbb{Z}$}{Z}}

The ring of integers is the structure
\[
  (\mathbb{Z}, +, \cdot, 0, 1).
\]
It provides the usual arithmetic operations and is also the setting for basic
number-theoretic algorithms such as greatest common divisor and primality
testing.

\subsection*{Greatest Common Divisor}

The greatest common divisor is computed using the Euclidean algorithm. For
integers $a,b$, it is characterized by
\[
\gcd(a,b) = \gcd(b, a \bmod b),
\]
with terminal cases
\[
\gcd(a,0)=|a|,
\qquad
\gcd(0,b)=|b|,
\qquad
\gcd(0,0)=0.
\]

\subsection*{Primality Testing}

Primality testing is based on the Miller--Rabin test. For integers in the
64-bit range, a fixed set of bases makes the test deterministic. Small factors
are first removed by trial division over
\[
\{2,3,5,7,11,13,17,19,23,29,31,37\},
\]
and then Miller--Rabin is run with that same set of bases. For
$n < 2^{64}$, this set of bases is sufficient to guarantee determinism
\cite{jaeschke1993miller, sorenson2015strong}.

\subsection*{Modular Rings and Polynomial Rings}

The modular ring $\mathbb{Z}/n\mathbb{Z}$ uses addition and multiplication
modulo $n$. It is a finite ring, and when $n$ is prime it is a field.

Given a coefficient ring $R$, the polynomial ring $R[X]$ consists of finite
formal sums
\[
a_0 + a_1X + \cdots + a_dX^d,
\qquad a_i \in R.
\]
Addition is coefficientwise, and multiplication is given by convolution:
\[
(fg)_k = \sum_{i+j=k} f_i g_j.
\]
